\documentclass[11pt,a4paper]{article}
\usepackage[utf8]{inputenc}
\usepackage[T1]{fontenc}
\usepackage[brazilian]{babel}
\usepackage{geometry}
\usepackage{enumitem}
\usepackage{listings}
\usepackage{xcolor}

\geometry{margin=2.2cm}

\lstset{
  basicstyle=\ttfamily\small,
  breaklines=true,
  frame=single,
  backgroundcolor=\color{gray!8}
}

\title{Verificação e Validação de Software \\ \large Trabalho 1: Defeitos Encontrados na Primeira Versão}
\author{Eduardo Alcaria \\ Arthur Pimentel \\ Ethan Soares}
\date{}

\begin{document}
\maketitle
\vspace{-2.5em}

\section*{Contexto}

A primeira versão da classe \texttt{CalculadoraTarifa} foi submetida ao conjunto
completo de 25 casos de teste unitários (arquivo
\texttt{CalculadoraTarifaTest.java}, técnica de particionamento e valor limite).
A execução resultou em 8 casos com falha, agrupados em 4 defeitos distintos,
descritos a seguir.

\section*{Defeito 1: cortesia com limite exclusivo}

\subsection*{Falha observada}
No limite exato de 20 minutos de permanência, a implementação cobrava a tarifa
fixa de R\$ 15,00 em vez de aplicar a cortesia (valor esperado R\$ 0,00).

\subsection*{Caso de teste que detectou}
CT02 (entrada 10/03/2026 08:00, saída 10/03/2026 08:20, sem VIP). Resultado
obtido: R\$ 15,00. Resultado esperado: R\$ 0,00.

\subsection*{Causa}
A condição de cortesia usava comparação estrita:
\begin{lstlisting}[language=Java]
if (minutos < MINUTOS_CORTESIA) {
    return 0.0;
}
\end{lstlisting}
O enunciado define a cortesia como ``20 minutos de cortesia'', ou seja, o
limite de 20 minutos deve estar incluído na faixa de valor zero.

\subsection*{Correção proposta}
Trocar o operador de comparação para incluir o limite:
\begin{lstlisting}[language=Java]
if (minutos <= MINUTOS_CORTESIA) {
    return 0.0;
}
\end{lstlisting}

\section*{Defeito 2: hora adicional cobrada apenas quando completa}

\subsection*{Falha observada}
Para permanências que ultrapassam a primeira hora sem completar um múltiplo
exato de 60 minutos, o valor cobrado ficava R\$ 5,00 abaixo do esperado, pois a
hora adicional iniciada não era cobrada.

\subsection*{Casos de teste que detectaram}
\begin{itemize}[nosep]
  \item CT06 (61 minutos): obtido R\$ 15,00, esperado R\$ 20,00.
  \item CT08 (121 minutos): obtido R\$ 20,00, esperado R\$ 25,00.
  \item CT12 (179 minutos): obtido R\$ 20,00, esperado R\$ 25,00.
  \item CT13 (359 minutos): obtido R\$ 35,00, esperado R\$ 40,00.
\end{itemize}

\subsection*{Causa}
O número de horas adicionais era calculado por divisão inteira, que trunca o
resultado em vez de arredondar para cima:
\begin{lstlisting}[language=Java]
long horasAdicionais = minutosExcedentes / MINUTOS_TARIFA_FIXA;
\end{lstlisting}
O enunciado especifica que o incremento ocorre ``a cada intervalo de 1 hora
(inclusive)'', ou seja, qualquer minuto iniciado dentro de uma nova hora já
deve gerar a cobrança daquele intervalo.

\subsection*{Correção proposta}
Arredondar o número de horas adicionais para cima:
\begin{lstlisting}[language=Java]
long horasAdicionais = (minutosExcedentes + MINUTOS_TARIFA_FIXA - 1) / MINUTOS_TARIFA_FIXA;
\end{lstlisting}

\section*{Defeito 3: pernoite não soma múltiplas diárias}

\subsection*{Falha observada}
Quando o veículo permanece por mais de uma pernoite, o sistema cobrava apenas
uma diária de R\$ 50,00, independentemente do número de noites em que o
veículo permaneceu estacionado.

\subsection*{Caso de teste que detectou}
CT15 (entrada 10/03/2026 20:00, saída 12/03/2026 08:00, duas pernoites
completas). Resultado obtido: R\$ 50,00. Resultado esperado: R\$ 100,00.

\subsection*{Causa}
A condição de pernoite era tratada como um valor booleano único, aplicando
sempre o mesmo valor fixo:
\begin{lstlisting}[language=Java]
boolean pernoite = saida.toLocalDate().isAfter(entrada.toLocalDate())
        && !saida.toLocalTime().isBefore(ABERTURA);
double valor = pernoite ? VALOR_PERNOITE : calcularPorPermanencia(minutos);
\end{lstlisting}
Isso não reflete o caso de um veículo que permanece por dois ou mais dias, no
qual o valor de pernoite deve ser cobrado uma vez para cada janela de 24 horas
excedida.

\subsection*{Correção proposta}
Contar quantas vezes a janela das 08:00 do dia seguinte é ultrapassada e
multiplicar pelo valor unitário de pernoite:
\begin{lstlisting}[language=Java]
long contarPernoites(LocalDateTime entrada, LocalDateTime saida) {
    LocalDateTime limite = entrada.toLocalDate().plusDays(1).atTime(ABERTURA);
    long pernoites = 0;
    while (!saida.isBefore(limite)) {
        pernoites++;
        limite = limite.plusDays(1);
    }
    return pernoites;
}
\end{lstlisting}

\section*{Defeito 4: janela de saída proibida com limites exclusivos}

\subsection*{Falha observada}
Uma saída registrada exatamente às 02:00 ou exatamente às 07:59 era aceita
pelo sistema, quando deveria ser rejeitada por cair dentro do intervalo em que
a saída de veículos não é permitida.

\subsection*{Casos de teste que detectaram}
\begin{itemize}[nosep]
  \item ES01 (saída às 02:00 do dia seguinte à entrada): esperava-se
        \texttt{IllegalArgumentException}, porém nenhuma exceção foi lançada.
  \item ES02 (saída às 07:59 do dia seguinte à entrada): esperava-se
        \texttt{IllegalArgumentException}, porém nenhuma exceção foi lançada.
\end{itemize}

\subsection*{Causa}
A verificação usava comparações estritas em ambas as extremidades do
intervalo:
\begin{lstlisting}[language=Java]
boolean dentroJanelaProibida = hora.isAfter(INICIO_SAIDA_PROIBIDA)
        && hora.isBefore(FIM_SAIDA_PROIBIDA);
\end{lstlisting}
O enunciado define o intervalo proibido como ``das 02:00 da manhã até as
07:59 da manhã'', portanto ambos os extremos devem ser considerados parte do
intervalo proibido.

\subsection*{Correção proposta}
Tornar a comparação inclusiva nos dois limites:
\begin{lstlisting}[language=Java]
boolean dentroJanelaProibida = !hora.isBefore(INICIO_SAIDA_PROIBIDA)
        && !hora.isAfter(FIM_SAIDA_PROIBIDA);
\end{lstlisting}

\section*{Observação adicional sobre o desenho dos testes}

Durante a análise, os casos ES01 e ES02 foram inicialmente escritos com a
saída no mesmo dia da entrada, o que fazia a validação de ordem cronológica
(saída deve ser posterior à entrada) mascarar o defeito real da janela de
saída proibida. Os casos foram reformulados posicionando a entrada no dia
anterior, isolando corretamente a condição sob teste.

\section*{Resultado após a correção}

Após aplicar as quatro correções descritas, a suíte completa de 25 casos de
teste passou a executar com sucesso (25 de 25), confirmando que os defeitos
identificados na primeira versão foram sanados sem introduzir regressões nos
demais casos.

\end{document}
